\scanpage{131}
\mspage{355}

the set of fixed \unsure{points} of $U_{0,3}$ under $\dot{\sigma}_\infty$ is reduced to a single element, which also implies that
\deleted{\illegible{1}}
\[
\Pi_{0,3}^{\sigma_\infty}=\{1\}.
\]

We shall \unsure{also} assume that these facts \unsure{extend} to the elements of $\widehat{\Pi}_{0,3}$ --- which will be the case if (in addition to what has been assumed about it with regard to $\rho$) we assume that we \unsure{still} have
\[
\tag{38}
\widehat{\Pi}_{0,3}.
\]

Note that here we have
\[
\tag{39}
G^{\sigma_\infty}=\mathcal T_{0,3}^{\sigma_\infty}\simeq\mathbb Z/2\mathbb Z\times\mathbb Z/2\mathbb Z,
\]
subgroup generated by $\dot{\sigma}_\infty,\unsure{\dot\tau|\tau}$, whose non\unsure{zero} elements are $\dot{\sigma}_\infty,\tau,\dot{\tau}'_\infty$ --- which \unsure{lift} to $\sigma_\infty,\tau_\infty,\tau'_\infty$ ; and we have
\[
\tag{40}
\mathcal E^{\sigma_\infty}=\mathcal G_{0,3}^{\sigma_\infty}\simeq\mathbb Z/2\mathbb Z\times\mathbb Z/2\mathbb Z,
\]
group generated by $\sigma_\infty,\tau_\infty$ (the other nontrivial element being $\tau'_\infty$). If then we write a solution $h\inserted{\in\pi}$ (of (35) in the form (36)), with $\beta\in\mathcal G_{0,3}$ only (not only $\beta\in\pi$), then the indeterminacy is in the group (40) of order (4) ; More precisely, if $\beta_0$ is the solution in $\pi$, then the other solutions are of the form
\[
\tag{41}
\beta=\beta_0\sigma_\infty^\varepsilon\tau_\infty^{\varepsilon'}
\qquad \varepsilon,\varepsilon'\in\mathbb Z/2\mathbb Z
\]

\margin{130}

\scanpage{132}

(with $\varepsilon,\varepsilon'$ uniquely determined). In order to have $\beta\in\mathcal G_{0,3}^{+}$, it is necessary and sufficient that $\varepsilon'=0$ !

\margin{131}

\scanpage{133}
\mspage{356}

\margin{These notations will be \unsure{revised} once again in \S\ 46}

\uline{44}\qquad \uline{Return to the notation.} \uline{Setting up equations} \uline{in} $\mathrm{Sl}(2,\widehat{\mathbb Z})$.

Let $u$ be an automorphism of $\widehat{\Pi}_{0,3}$ \inserted{commuting up to inner automorphisms : $\unsure{\mathcal G_3}$\deleted{\illegible{1}}} --- or what amounts to the same, an automorphism of $\widehat{\mathcal G}_{0,3}^{+}$ which induces the identity in $\mathcal T_{0,3}^{+}=\unsure{\mathcal G_3}$. Then $u$ is known when we know its actions on the two generators $\sigma_\infty,\rho$ of $\widehat{\mathcal G}_{0,3}^{+}$, which will resp. be elements of $\widehat{\mathcal G}_{0,3}^{+}$ over $\dot{\sigma}_\infty$ and $\dot\rho$,

namely
\[
\tag{1}
\boxed{u(\sigma_\infty)=h\sigma_\infty,\qquad u(\rho)=g\rho}
\qquad (h,g\in\widehat{\Pi}_{0,3})
\]
where $h,g$ \uline{are subject to the conditions}
\[
\tag{2}
\boxed{h\sigma_\infty(h)=1\qquad g\rho(g)\rho^2(g)=1},
\]
expressing respectively the relations
\[
\tag{3}
\boxed{(h\sigma_\infty)^2=(g\rho)^3=1}.
\]

(There is in addition the \unsure{preliminary} condition to check --- namely that (3) are defining relations for $\widehat{\mathcal G}_{0,3}^{+}$ \dots).

\deleted{\illegible{1}}\inserted{For the moment}, there is no need to impose any relation between $g$ and $h$ --- except the relations (2) on each of these elements --- since $\widehat{\mathcal G}_{\unsure{0,3}}^{+}$ is the group with defining relations
\[
\sigma_\infty^2=\rho^3=1.
\]

\margin{132}

\scanpage{134}

The elements $g$ and $h$, thus taken \unsure{as} automorphisms of $\widehat{\Pi}_{0,3}$, express the commutation relations of\deleted{s} $u$ with \deleted{$\sigma_\infty$} $\rho$ and $\sigma_\infty$ respectively :
\[
\tag{4}
[u,\rho]=\underbrace{u\rho u^{-1}}_{u(\rho)}\rho^{-1}=\operatorname{int}(g),
\qquad
[u,\sigma_\infty]=u\sigma_\infty u^{-1}\sigma_\infty=\operatorname{int}(h)
\]
--- it is a remarkable thing that knowledge of the elements $g,h$ of $\pi$ satisfying (4), already \uline{\unsure{implies|imply}} knowledge of $u$ (\unsure{this|which} \unsure{to the} fact that $\sigma_\infty,\rho$ generate $\mathcal G_{0,3}^{+}$ \dots).

Let us now write the compatibility of $u$ with the loop structure, \deleted{\illegible{3}}\inserted{We had it} \deleted{\unsure{formulated}}\inserted{\unsure{expressed|expressed}} previously by a relation of the form
\[
u(\varepsilon_0)=\operatorname{int}(g_0)(\varepsilon_0^p)
\quad\text{or}\quad
u(\ell_0)=\operatorname{int}(g_0)\ell_0^p,
\]
for
\[
p\in\widehat{\mathbb Z}^{*},\qquad g_0\in\widehat{\Pi}_{0,3}
\]
suitably chosen. The notation $g_0$ is not all that happy, since $g$ and $g_0$ play quite dissimilar roles. The subscript $0$ in $g_0$ was intended to allow us to write, more generally
\[
\deleted{\unsure{g_i}}\quad
u(\ell_i)=\operatorname{int}(g_i)\ell_i^p
\qquad i\in\{0,1,\infty\}.
\]

\margin{133}

\scanpage{135}
\mspage{357}

I prefer to use the letter $f_0$ --- with the option of dropping the subscript if no confusion is possible. The letter $p$ for a multiplier is not \inserted{a choice} that is very happy, for later on \deleted{\unsure{$p$}} we shall \inserted{\unsure{surely}} need this letter to denote a characteristic --- I prefer the letter $\mu$ (Greek initial of ``multiplier''), and the loop relation \inserted{\unsure{would}} therefore take the form
\[
\deleted{\unsure{(5)}}\qquad
u(\varepsilon_0)=\operatorname{int}(f_0)\varepsilon_0^\mu.
\qquad
\deleted{\unsure{\text{or}}\quad u(\ell_0)=\operatorname{int}(f_0)}
\]

In fact, I prefer to replace $f_0$ in it by $f_0^{-1}$, to obtain later on formulas in which the symmetry of the roles played by $\rho$ and by $\sigma_\infty$ \unsure{appears} more clearly. --- in particular I would like the formula defining $h_0$ in terms of $g_0$ (\S\ 43, formula (65,)
\[
h_0=g_0^{-1}\sigma_\infty(g_0)
\]
to become instead
\[
h_0=g_0\sigma_\infty(g_0)^{-1},
\]
in harmony with
\[
\xi=\lambda\rho(\lambda)^{-1}\quad\dots
\]

Thus the ``\unsure{loop|loop} relation'' takes the form
\[
\tag{5}
u(\varepsilon_0)=\operatorname{int}(f_0^{-1})\varepsilon_0^\mu
\qquad f_0\in\widehat{\Pi}_{0,3},\ \mu\in\widehat{\mathbb Z}^{*},
\]
\deleted{\unsure{rather} \unsure{we} immediately deduce}
\[
\tag{6}
u(\ell_0)=\operatorname{int}(f_0^{-1})\ell_0^\mu,
\]
from which we immediately deduce

\margin{134}

\scanpage{136}

\[
\tag{7}
\left\{
\begin{aligned}
u(\varepsilon_1)&=\operatorname{int}(f_1^{-1})(\varepsilon_1^\mu),
&u(\ell_1)&=\operatorname{int}(f_1^{-1})(\ell_1^\mu),\\
u(\varepsilon_\infty)&=\operatorname{int}(f_\infty^{-1})(\varepsilon_\infty^\mu),
&u(\ell_\infty)&=\operatorname{int}(f_\infty^{-1})(\ell_\infty^\mu),
\end{aligned}
\right.
\]
\unsure{where|or}
\[
\tag{8}
f_1=\rho(f_0)g^{-1},\qquad
f_\infty=\rho(f_1)g^{-1}=\rho^2(f_0)\rho(g)^{-1}g^{-1}.
\]

In what follows, we shall simply write $f$ instead of $f_0$, and we are going to transform equation (5) \deleted{\illegible{1}} into
\[
\tag{8bis}
\operatorname{int}(f)u(\varepsilon_0)=\varepsilon_0^\mu,
\]
by making the left-hand side explicit
\[
\begin{aligned}
\operatorname{int}(f)u(\varepsilon_0)
&=f\,u(\sigma_\infty\rho)f^{-1}
=f\underbrace{u(\sigma_\infty)}_{h\sigma_\infty}\underbrace{u(\rho)}_{g\rho}f^{-1}
=fh\sigma_\infty g\rho f^{-1}\\
&=fh\sigma_\infty g\sigma_\infty^{-1}\underbrace{\sigma_\infty\rho}_{\varepsilon_0}f^{-1}
=fh\sigma_\infty(g)\underbrace{\varepsilon_0f^{-1}\varepsilon_0^{-1}}\varepsilon_0
\end{aligned}
\]
i.e.
\[
\tag{9}
\operatorname{int}(f)u(\varepsilon_0)
=\bigl[fh\sigma_\infty(g)(\sigma_\infty\rho)(f)^{-1}\bigr]\varepsilon_0
=\bigl[f,h\sigma_\infty(g\rho(f)^{-1}\bigr]\varepsilon_0
\]
where the \unsure{four|three} factors $f,h,\sigma_\infty(g\rho(f)^{-1}$ in $[\ ]$ are $\in\widehat{\Pi}_{0,3}$. Thus \unsure{(5)|(9)} or (8), takes the form
\[
fh\sigma_\infty(g\rho(f)^{-1})=\varepsilon_0^{\mu-1}\,\deleted{\illegible{1}}
\]
which gives, setting
\[
\tag{10}
\mu-1=2\nu\qquad(\nu\in\widehat{\mathbb Z})
\]
(so we write $\nu$ for what was denoted by $\gamma$ previously) we find the relation
\[
fh\sigma_\infty\bigl(g\rho(f)^{-1}\bigr)=\ell_0^\nu
\]
or again, transforming by $\sigma_\infty$
\[
\sigma_\infty(\unsure{f})\sigma_\infty(h)g\rho(f)^{-1}=\ell_1^\nu
\]
; that is

\margin{135}

\scanpage{137}
\mspage{358}

\[
\boxed{\sigma_\infty(h).g=\sigma_\infty(f)^{-1}\ell_1^\nu\rho(f)}.
\qquad (11)
\]

This relation, which relates $h$ and $g$ via $f$, can rightly be called the ``loop relation''.
We see, from this relation, that \deleted{\illegible{1}} not only does knowledge \inserted{of $u$ i.e.} \unsure{of the|of} $h$ and $g$ determine $f$ and $\mu$ i.e. $f$ and $\nu$, but also conversely, once $f$ and $\nu$ are known, $g$ and $h$ determine each other via (11).

\margin{$f$ determined modulo $L_0$ only $f\mapsto\ell_0^n f$ $(n\in\widehat{\mathbb Z})$.}

Note that for $u$ corresponding to : \deleted{\unsure{$f$|$\xi$}}, $g,h,f$, if $\gamma\in\widehat{\Pi}_{0,3}$, the \unsure{automorphism} \inserted{$u'=$} $\operatorname{int}(\gamma)\circ u$ corresponds to
\[
g'=\gamma g\rho(\gamma^{-1}),\quad
h'=\gamma h\sigma_\infty(\gamma)^{-1},\quad
f'=f\gamma^{-1}.
\qquad (12)
\]

\margin{This holds more \unsure{generally} if $\gamma\in\unsure{\widetilde{\mathcal G}_{0,3}|\mathcal G_{0,3}}$, and even for $\gamma\in\unsure{\widehat{\widetilde{\mathcal G}}_{0,3}|\widehat{\mathcal G}_{0,3}}$.}

The equations (2) ``in $h$'' and ``in $g$'' are solved parametrically in a unique way by
\[
g=\alpha\rho(\alpha)^{-1},\quad
h=\beta\sigma_\infty(\beta)^{-1}
\qquad (13)
\]
where
\[
\left\{\begin{array}{l}
\alpha\in\widehat{\widetilde{\Pi}}_{0,3}=\widehat{\Pi}_{0,3}.\{1,\tau_\infty\}\\
\beta\in\widehat{\Pi}_{0,3}
\end{array}\right.
\qquad (14)
\]
to be precise, giving $g,h$ satisfying (2) is equivalent resp. to giving $\alpha,\beta$ in $\widehat{\widetilde{\Pi}}_{0,3}$ and $\widehat{\Pi}_{0,3}$, via the formulas (13).

\margin{136}

\scanpage{138}

\inserted{The elements} $g$ and $h$ being thus made explicit in terms of $\alpha$ and $\beta$, the \inserted{loop} relation (11) can be considered as a relation between $\alpha,\beta,f,\nu$
\[
\sigma_\infty(\beta\sigma_\infty(\beta)^{-1})(\alpha\rho(\alpha)^{-1})
=\sigma_\infty(f)^{-1}\ell_1^\nu\rho(f),
\qquad (15)
\]
where \deleted{$\nu,f$ \unsure{are}} $\nu$ is uniquely determined by $\alpha,\beta$, and $f$ is determined in $L_0\backslash\widehat{\Pi}_{0,3}$.
(We shall note that the right-hand side of (15) does not change, on changing $f$ to $\ell_0^n f$, with $n\in\widehat{\mathbb Z}$).

The main asymmetry, it seems to me, between $\sigma_\infty$ and $\rho$ in this theory, is that \deleted{\illegible{1}} we always have $\beta\in\widehat{\Pi}_{0,3}$, whereas $\alpha$ is not necessarily in $\widehat{\Pi}_{0,3}$ -- it is, \unsure{as we have seen}, iff $\nu\equiv0\ (3)$ i.e. $\mu\equiv1\ (3)$.

\deleted{So if we use \illegible{1} $\operatorname{int}(\beta^{-1})$ \illegible{1}}
The automorphism
\[
u'=\operatorname{int}(\beta^{-1})u
\]
will therefore be \unsure{congruent} to $u$ modulo \unsure{inner} automorphisms, and will correspond to
\[
\left\{\begin{array}{l}
g'=\beta^{-1}g\rho(\beta)=(\beta^{-1}\alpha)\rho(\beta^{-1}\alpha)^{-1}\\
h'=\beta^{-1}h\sigma_\infty(\beta)=1\\
f'=f\beta
\end{array}\right.
\qquad (16)
\]
-- thus $u$ is \deleted{conjugate} \inserted{congruent modulo} inner automorphisms

\margin{137}

\scanpage{139}
\mspage{359}

to a unique $u'$ such that $u'(\sigma_\infty)=\sigma_\infty$ i.e. a $u'$ which commutes with $\sigma_\infty$. But in \unsure{the symmetric notation|the symmetric notations} we shall take
\[
u''=\operatorname{int}(\alpha^{-1})u
\]
corresponding to
\[
\left\{\begin{array}{l}
g''=\alpha^{-1}g\rho(\alpha)=1\\
h''=\alpha^{-1}h\sigma_\infty(\alpha)=(\alpha^{-1}\beta)\sigma_\infty(\alpha^{-1}\beta)^{-1}\\
f''=f\alpha
\end{array}\right.
\qquad (17)
\]
but $u''$ will be \inserted{only $\widehat{\widetilde{\Pi}}_{0,3}$-}conjugate \deleted{\unsure{by}} of $\unsure{u''|u}$ -- it will be \inserted{$\widehat{\Pi}_{0,3}$-}conjugate to an \deleted{element of} \inserted{\unsure{automorphism}} which commutes with $\rho$ iff $g$ is $\Pi_{0,3}$-$\rho$-conjugate to $1$ i.e. iff $\alpha\in\widehat{\Pi}_{0,3}$ i.e. iff $\mu\equiv1\ (3)$ or again $\nu\equiv0\ (3)$.

\deleted{As for the} The calculations of the preceding \S\S\ concerned the case where $h=1$, i.e. $u$ commutes with $\sigma_\infty$ -- so \illegible{1} there remain the quantities $g$ and $f$, \deleted{we see moreover} \inserted{plus the multiplier $\mu$}
\deleted{\inserted{the multiplier} for $\mu$, \unsure{or for} $\nu=\frac{\mu-1}{2}$, is determined unambiguously by}
-- once $\mu$ is known i.e. $\nu$ known, $g$ and $f$ determine one another by (11) ($f$ remaining an element of $L_0\backslash\widehat{\Pi}_{0,3}$).

\margin{We shall see that $\mu$ is already determined by $f$, \unsure{at least modulo} sign, cf. below.}

For fixed $f$, the equation in $g$ can be interpreted as an equation in $f$, or better, introducing
\[
\varphi=f\sigma_\infty(f)^{-1}
\qquad (12)
\]

\margin{This notation is \unsure{useful|valid} only if $\beta=1$ \illegible{2}}
\margin{it is essentially the element denoted $h_0$ previously}
\margin{138}

\scanpage{140}

an ``equation in $\varphi$'' --
\[
(\ell_0^\nu\varphi)\rho(\ell_0^\nu\varphi)\rho^2(\ell_0^\nu\varphi)=1,
\qquad (13)
\]
or again, putting
\[
\boxed{\xi=\ell_0^\nu\varphi\quad\text{i.e. }\varphi=\ell_0^{-\nu}\xi}
\qquad (14)
\]
we can interpret it as ``equation in $\xi$''
\[
\xi\rho(\xi)\rho^2(\xi)=1.
\qquad (15)
\]
Thus, the loop automorphisms of $\widehat{\Pi}_{0,3}$ which commute with $\sigma_\infty$, \inserted{with multiplier $\mu=2\nu+1$,} correspond to the pairs of solutions $(\xi,\varphi)$ of the equations
\[
\boxed{\varphi\sigma_\infty(\varphi)=1\qquad
\xi\rho(\xi)\rho^2(\xi)=1}
\qquad (16)
\]

\margin{We \unsure{shall find|find} that \unsure{both} $\xi$ and $\varphi$ already \unsure{determine} $\mu$, hence $\nu$.}

satisfying in addition the compatibility condition (14) (particularly simple \deleted{\unsure{(17)}} if $\mu=1$ i.e. \deleted{$\mu=$} $\nu=0$, since then $\varphi=\xi$ !). We recover $g$ (hence $u$) \deleted{by the} as a function
\deleted{Conversely, if we \illegible{3}}
of $(\varphi,\xi)$ by the loop formula (11), which here takes the form
\[
g=\sigma_\infty(f)^{-1}\ell_1^\nu\rho(f)
=(\ell_0^\nu f)^{-1}\bigl(\underbrace{\ell_0^\nu\varphi}_{\xi}\bigr)\rho(\ell_0^\nu f)
\qquad (17)
\]
$f\in\widehat{\Pi}_{0,3}$ being uniquely determined in terms of $\varphi$ by formula (12).

Symmetrically, we are going to look at the loop automorphisms $u$ of $\Pi_{0,3}$ which com-

\margin{139}

\scanpage{141}
\mspage{360}

mute with $\rho$, i.e. those for which we have $g=1$, hence described in terms of $h,f$,
\[
u(\sigma_\infty)=h\sigma_\infty,\quad u(\rho)=1_{\deleted{\unsure{\rho}}}\qquad (18)
\]
\[
u(\varepsilon_0)=\operatorname{int}(f^{-1})(\varepsilon_0^\mu)
\qquad (19)
\]
this last relation taking the form (11)
\[
\sigma_\infty(h)=\sigma_\infty(f)^{-1}\ell_1^\nu\rho(f),
\qquad (20)
\]
which can also be written, applying $\sigma_\infty$
\[
h=f^{-1}.\ell_0^\nu.(\sigma_\infty\rho)(f)\unsure{\underset{\cdot}{=}|=}
\qquad (21)
\]
``The equation in $h$'' (2) \unsure{the}
\[
h\sigma_\infty(h)=1
\qquad \unsure{(*)}
\]
is thus written as an ``equation in $f$'', \deleted{which} \inserted{with}
\[
h\sigma_\infty(h)=f^{-1}\ell_0^\nu
\underbrace{(\sigma_\infty\rho)(f)\sigma_\infty(f)^{-1}}_{\sigma_\infty(\underbrace{\rho(f)f^{-1}}_{\psi^{-1}})}
\ell_1^\nu\rho(f)\deleted{\illegible{1}}
\]
so \unsure{this} is written, in terms of $\psi=f\rho(f)^{-1}$
\[
\underbrace{\ell_0^\nu\sigma_\infty(\psi^{-1})}\ell_1^\nu
\underbrace{\rho(f)f^{-1}}_{\psi^{-1}}=1
\]
i.e. (passing to the inverse)
\[
\psi\ell_1^{-\nu}\underbrace{\sigma_\infty(\psi)\ell_0^{-\nu}}_{\sigma_\infty(\psi\ell_1^{-\nu})}=1
\quad\text{i.e.}\quad
\underbrace{\ell_0^{-\nu}\psi}_{\eta}
\underbrace{\ell_1^{-\nu}\sigma_\infty(\psi)}_{\sigma_\infty(\underbrace{\ell_0^{-\nu}\psi}_{\eta})}=1.
\]
\deleted{i.e.}
Thus, putting
\[
\psi\overset{\mathrm{dfn}}{=}f\rho(f)^{-1}
\qquad (22)
\]
\[
\boxed{\eta=\ell_0^{-\nu}\psi\quad\text{i.e. }\psi=\ell_0^\nu\eta}
\qquad (23)
\]
, we find that the loop automorphisms $u$ of $\widehat{\Pi}_{0,3}$ which commute with $\rho$ correspond to the triples $(\eta,\psi,\nu)$, $\eta,\psi\in\widehat{\Pi}_{0,3}$, $\nu\in\widehat{\mathbb Z}$ (with $2\nu+1\in\widehat{\mathbb Z}^{*}$), satisfying

\margin{140}

\scanpage{142}

\[
\boxed{\psi\rho(\psi)\rho^2(\psi)=1,\qquad\eta\sigma_\infty(\eta)=1},
\qquad (24)
\]
\margin{We \unsure{shall find} \illegible{1} \unsure{again}, \unsure{by} \unsure{both} $\psi$ and $\eta$}

and in addition the compatibility relation (23) -- particularly simple if $\mu=1$ i.e. if $\nu=0$, for then it means $\psi=\eta$ ! But in addition to these relations, we must assume that the solution $\psi$ of equation (24) ``in $\psi$'' \unsure{belongs|belong} to the principal family, i.e. that, when we write $\psi$ in parametric form
\[
\psi=f\rho(f)^{-1}
\qquad (22)
\]
with $f\in\widehat{\widetilde{\Pi}}_{0,3}=\widehat{\Pi}_{0,3}.\{1,\tau_\infty\}$, we indeed have $f\in\widehat{\Pi}_{0,3}$ (hence that the image of $f$ in $\widehat{\widetilde{\Pi}}_{0,3}/\widehat{\Pi}_{0,3}\simeq\{1,\tau_\infty\}\simeq\mathbb Z/2$ is zero\dots).

Having thus determined $f$ as a \uline{function} of $\psi$, \unsure{determine|determines} $h$ by (21), whence $u$ \deleted{\illegible{3}}.

\uline{Remarks} 1) If this additional condition on $\psi$ is not satisfied, the \deleted{\unsure{determination}} formula (21) still makes sense and defines $h\in\widehat{\Pi}_{0,3}$, satisfying ``the equation in $h$'' $h\sigma_\infty(h)=1$ hence defining an automorphism $u$ of $\widehat{\Pi}_{0,3}$, commuting with $\rho$.

We shall then have
\[
u(\ell_0)=\operatorname{int}(f^{-1})(\ell_0^\mu)
\]
or again, putting
\[
f=\inserted{\tau_\infty}f_0\deleted{\tau_\infty}\qquad\text{with }f_0\in\widehat{\Pi}_{0,3}
\]
\[
u(\ell_0)=\operatorname{int}(f_0^{-1})(\tau_\infty(\ell_0^\mu))=\operatorname{int}(f_0^{-1})(\ell_0^{-\mu})
\]

\margin{141}

\scanpage{143}
\mspage{361}

i.e. $u$ is a loop automorphism, with multiplier $-\mu=-(2\nu+1)$.

2) Let us note that the indeterminacy of $f$ (for fixed $u$, i.e. fixed $g,h$) lies in
\[
 f\longmapsto f'=\ell_0^n f\;\deleted{\unsure{f}}
 \qquad(n\in\widehat{\mathbb Z})
 \tag{25}
\]
which gives
\[
 \left\{\begin{aligned}
 \varphi&\longmapsto\varphi'=\ell_0^n\varphi\ell_1^{-n},&
 \psi&\longmapsto\psi'=\ell_0^n\psi\ell_1^{-n},\\
 \xi&\longmapsto\xi'=\ell_0^n\xi\ell_1^{-n},&
 \eta&\longmapsto\eta'=\ell_0^n\eta\ell_1^{-n},
 \end{aligned}\right.
 \tag{26}
\]
\margin{\uline{N.B.} We similarly have
$\alpha_0\longmapsto\alpha_0'=\ell_0^n\alpha_0$
$\beta_0\longmapsto\beta_0'=\ell_0^n\beta_0$.}

(N.B. we have $\ell_1=\sigma_\infty(\ell_0)=\rho(\ell_0)$, so $\xi'$ and $\psi'$ are the $\rho$-transforms of $\xi,\psi$ by $\unsure{\ell_0^{-n}|\ell_0^n}$, and $\varphi',\eta'$ the $\sigma_\infty$-transforms of $\varphi,\eta$ by $\unsure{\ell_0^{-n}|\ell_0^n}$ --- in particular the equations (16) (in $\varphi$ and in $\xi$) and (24) (in $\psi$ and in $\eta$) are indeed invariant under these transformations, as are moreover the compatibility relations (14) resp. (23).

3) We have the following formulas \inserted{(rewriting of (11))}
\[
 \sigma_\infty(h)g
 = (\ell_0^\nu f)^{-1}\xi\rho(\ell_0^\nu f)
 = \sigma_\infty(f^{-1}\ell_0^\nu)\eta(f^{-1}\ell_0^\nu)^{-1}
 \tag{27}
\]
\deleted{(case $h=1$)} --- expressing $\inserted{\sigma_\infty(h)}g$ as the $\rho$-conjugate of $\xi$ by $(\ell_0^\nu f)^{-1}$, and symmetrically \deleted{in the case $g=1$,} as the $\sigma_\infty$-conjugate of $\eta$ by $\sigma_\infty(f^{-1}\ell_0^\nu)$. If $h=1$, it is the first expression that is interesting (implying the equation in $\xi$), for $g=1$, it is the second, which we can also rewrite \inserted{(applying $\sigma_\infty$ to the 2 sides)}
\[
 h=(f^{-1}\ell_0^\nu)(\eta^{-1})\sigma_\infty(f^{-1}\ell_0^\nu)^{-1}
 \tag{28}
\]
\margin{142}

\scanpage{144}

expressing $h$ as $\sigma_\infty$-transform of $\eta^{-1}$ by $f^{-1}\ell_0^\nu$. If we want to express it as \inserted{$\sigma_\infty$-}transform of $\eta$ (instead of $\inserted{\sigma_\infty(\eta)=}\eta^{-1}$), we write
\[
 \deleted{\unsure{\eta^{-1}=\eta^{-1}\eta\sigma_\infty(\eta^{-1})}}
\]
\[
 \eta^{-1}=\sigma_\infty(\eta)=\eta^{-1}.\eta.\sigma_\infty(\eta)
\]
whence
\[
 h=(f^{-1}\ell_0^\nu\eta^{-1})\eta\sigma_\infty(f^{-1}\ell_0^\nu\eta^{-1})^{-1}.
 \tag{28bis}
\]

These formulas allow us to relate, in the case $h=1$ the solutions of the \deleted{\unsure{two}} equations
\[
 g=\alpha\rho(\alpha)^{-1}
 \qquad\text{and}\qquad
 \xi=\alpha_0\rho(\alpha_0)^{-1}
 \tag{29}
\]
by
\[
 \alpha=(\ell_0^\nu f)^{-1}\alpha_0=f^{-1}\ell_0^{-\nu}\alpha_0,
 \tag{30}
\]
and in the case where $g=1$, the solutions of the two equations
\[
 h=\beta\sigma_\infty(\beta)^{-1},\qquad
 \eta=\beta_0\sigma_\infty(\beta_0)^{-1}
 \tag{31}
\]
by
\[
 \beta=(f^{-1}\ell_0^\nu\eta^{-1})\beta_0.
 \tag{32}
\]

The asymmetry between these two formulas ((32) being a little more complicated than (30) by the extra factor $\eta^{-1}$ in the coefficient of $\beta_0$) comes from the fact that $f$ was defined via the loop \unsure{formula|formulas} (6), first form, where $\ell_0$ is defined as the square of $\sigma_\infty\rho=\varepsilon_0$, and not of $\rho\sigma_\infty=\varepsilon_1$
\margin{143}

\scanpage{145}
\mspage{362}

If we now identify $\widehat{\mathcal G}_{0,3}$ with $\mathrm{Gl}(2,\mathbb Z)^{\prime\wedge}$, and look at the \deleted{two} images of the elements $g,h,f$ and so on in the quotient
\[
 \mathrm{Gl}(2,\widehat{\mathbb Z})^{\pm1\prime}
 =\mathrm{Gl}(2,\widehat{\mathbb Z})^{\pm1}/\pm1
\]
--- which we denote by the same letters, by abuse of language --- \inserted{we} are led to calculate with elements of $\mathrm{Gl}(2,\widehat{\mathbb Z})$, rather than of $\mathrm{Gl}(2,\widehat{\mathbb Z})'$. We \unsure{will choose} $\alpha,\beta$ in $\widehat{\Pi}_{0,3}$ and $\widehat{\widetilde{\Pi}}_{0,3}$ in accordance with (13), hence
\[
 \alpha,\beta\in\mathrm{Gl}(2,\widehat{\mathbb Z})^{\pm1}
 \quad\text{\uline{modulo sign}},
 \tag{33}
\]
but the corresponding elements
\[
 g=\alpha\rho(\alpha)^{-1},\quad
 h=\beta\sigma_\infty(\beta)^{-1}
 \quad\in\mathrm{Sl}(2,\widehat{\mathbb Z})^{\deleted{\unsure{\pm1}}}
 \tag{34}
\]
are determined in $\mathrm{Gl}(2,\widehat{\mathbb Z})^{\deleted{\unsure{\pm1}}}$ \uline{without sign ambiguity}. Likewise
\[
 \overset{\inserted{\text{(case }h=1\text{)}}}{\alpha_0}
 \;\unsure{\text{and}|,}\;
 \overset{\inserted{\text{(case }g=1\text{)}}}{\beta_0}
 \in\mathrm{Gl}(2,\widehat{\mathbb Z})^{\pm1}
 \quad\text{modulo sign},
 \tag{35}
\]
and then
\[
 \begin{aligned}
 \xi&=\alpha_0\rho(\alpha_0)^{-1}
 \inserted{\quad\text{(case }h=1\text{)}},\quad
 \eta=\beta_0\sigma_\infty(\beta_0)^{-1}
 \inserted{\quad\text{(case }g=1\text{)}}
 \in\mathrm{Sl}(2,\widehat{\mathbb Z})
 \end{aligned}
 \tag{36}
\]
are determined without sign ambiguity. We have
\deleted{\unsure{This done|This fixed} \unsure{(34)|(37)}, the elements}
\deleted{the relations}
\[
 \deleted{\unsure{\varphi=} }f\in\mathrm{Sl}(2,\widehat{\mathbb Z})
 \quad\text{modulo \uline{sign}}
 \quad\text{(for }f\in\widehat{\Pi}_{0,3}\text{ fixed)}
 \tag{37}
\]
and we deduce from this
\[
 \varphi=f\sigma_\infty(f)^{-1},\quad
 \psi=f\rho(f)^{-1}\in\mathrm{Sl}(2,\widehat{\mathbb Z})
 \tag{38}
\]
without sign ambiguity.

\deleted{The relations} In the choice of the \deleted{\unsure{repr}} representatives

\deleted{\unsure{(39)}\ \illegible{1}\ $\unsure{\varphi|\xi}=$}
\margin{144}

\scanpage{146}

of $\varepsilon_0\unsure{\in}\mathrm{Sl}(2,\widehat{\mathbb Z})'$ in $\mathrm{Sl}(2,\widehat{\mathbb Z})$ there is sign ambiguity
\[
 \varepsilon_0=\pm\begin{pmatrix}1&-1\\0&1\end{pmatrix},
\]
\margin{though the sign $-$ would \unsure{hardly} be a \unsure{reasonable} choice --- one likes to see a unipotent element \deleted{\illegible{1}} generating a 1-parameter \unsure{subgroup}\ldots}
but for
\[
 \ell_0=\varepsilon_0^2=\begin{pmatrix}1&-2\\0&1\end{pmatrix}
 \tag{39}
\]
there is no ambiguity of \uline{sign}. This being so, the relations
\[
 \xi=\ell_0^\nu\varphi\qquad\text{(case }h=1\text{)}
 \tag{40}
\]
and
\[
 \eta=\ell_0^{-\nu}\psi\qquad\text{(case }g=1\text{)}
 \tag{41}
\]
are they true without a correcting \uline{sign} \unsure{??|?!}

The case $h=1$ corresponds to the data (in $\widehat{\widetilde{\Pi}}_{0,3}$) of $\xi$ and $\deleted{\unsure{\varphi}}f$, \inserted{and $\nu$,} satisfying
\[
 \xi\rho(\xi)\rho^2(\xi)=1,\qquad
 \xi=\ell_0^\nu
 \bigl(\underbrace{f\sigma_\infty(f)^{-1}}_{\varphi}\bigr),
\]
\deleted{and the equation in $f$ \unsure{is}}
or again to the data
\deleted{$\alpha_0\in\widehat{\widetilde{\Pi}}_{0,3},\ f\in\widehat{\Pi}_{0,3},$}
\[
 \alpha_0\in\widehat{\widetilde{\Pi}}_{0,3},\quad
 f\in\widehat{\Pi}_{0,3},\quad\nu\in\widehat{\mathbb Z}
\]
satisfying the relations
\[
 \underbrace{\alpha_0\rho(\alpha_0)}_{\xi}
 =\ell_0^\nu\underbrace{f\sigma_\infty(f)^{-1}}_{\varphi}
 \tag{42}
\]
in terms of these data we define
\[
 \left\{\begin{aligned}
 \xi&=\alpha_0\rho(\alpha_0)\\
 \varphi&=f\sigma_\infty(f)^{-1}\\
 g&=\sigma_\infty(f)^{-1}\ell_1^\nu\rho(f)
 =(\ell_0^\nu f)^{-1}(\ell_0^\nu\varphi)\rho(\ell_0^\nu f)
 \end{aligned}\right.
 \tag{43}
\]
This recalled, we lift $\alpha_0,f$ to $\mathrm{Gl}(2,\widehat{\mathbb Z})^{\pm1}$,
\margin{145}

\scanpage{147}
\mspage{363}

with sign ambiguity, \deleted{\unsure{but we}} and \deleted{\illegible{2}}\inserted{$\xi$ and $\varphi$} are thus determined in $\unsure{\mathrm{Sl}|\mathrm{Gl}}(2,\widehat{\mathbb Z})$, independently of the choices of signs, \deleted{\illegible{2}} and likewise $g$.

We \unsure{will have|have} therefore in $\mathrm{Sl}(2,\widehat{\mathbb Z})$
\[
 \deleted{\xi=\alpha_0\rho(\alpha_0)=\xi\,\unsure{\varepsilon\ell_0^\nu}}
\]
\[
 \xi=\varepsilon\ell_0^\nu\varphi
 \qquad\text{i.e.}\qquad
 \alpha_0\rho(\alpha_0)=\varepsilon\ell_0^\nu f\sigma_\infty(f)^{-1},
 \tag{44}
\]
\inserted{\uline{in $\mathrm{Sl}(2,\widehat{\mathbb Z})$},}
with a sign $\varepsilon\in\{\pm1\}$ that is well determined, and the question that arises is whether we \unsure{always|always} have $\varepsilon=+1$.

\deleted{\unsure{Defining} $f,\unsure{\rho_0|\rho}$}
\inserted{Let us note that the formula}
\[
 \rho(\ell_0)=\ell_1
\]
is true in $\mathrm{Sl}(2,\widehat{\mathbb Z})$ (without a correcting \uline{sign}) so the \inserted{last} expression for $g$ in (43) is valid in $\mathrm{Sl}(2,\widehat{\mathbb Z})$ (without a \uline{sign}), which shows that $g\in\mathrm{Sl}(2,\widehat{\mathbb Z})$, \deleted{\unsure{is not}} is a $\rho$-transform of $\ell_0^\nu\varphi$, \deleted{\unsure{therefore}} the latter is a $\rho$-transform of $g$.
\deleted{\illegible{1} $g$, In fact (\illegible{3} $g$) the relation}
\inserted{We have}
\[
 g\rho(g)\rho^2(g)=\varepsilon_0.1
 \qquad\varepsilon\in\{\pm1\},
\]
\[
 \deleted{(\ell_0^\nu\varphi)\rho(\ell_0^\nu\varphi).}
\]
and \deleted{likewise} we deduce from this, if $g'=\alpha g\rho(\alpha)^{-1}$
\deleted{\unsure{$\rho^2\in$}} ($\alpha\in\mathrm{Gl}(2,\widehat{\mathbb Z})$), that we have
\[
 g'\rho(g')\rho^2(g')=\varepsilon_0.1
\]
hence in particular for $g'=\ell_0^\nu\varphi\inserted{=\varepsilon\xi\quad\text{(44)}}$ \deleted{hence}, with $\varepsilon\in\{\pm1\}$ \deleted{\illegible{2}} now $\xi$ given by (42) satisfies
\[
 \xi\rho(\xi)\rho^2(\xi)=1
\]
(\uline{N.B.} we do not have $\rho^3=1$ in $\mathrm{Sl}(2,\widehat{\mathbb Z})$, but $\rho^3=-1$, which implies that the inner automorphism \deleted{\unsure{that it}} \unsure{defined|defines} by $\rho$ is of order 3\ldots), hence
\[
 \varepsilon_0=\varepsilon,
\]
\margin{146}

\scanpage{148}

i.e. the question of the validity of the formula
\[
 \xi=\ell_0^\nu\varphi
\]
in $\mathrm{Sl}(2,\widehat{\mathbb Z})$ is equivalent to that of the validity of the formula
\[
 g\rho(g)\rho^2(g)=1,
\]
which we leave open for the moment. Let us note for reference the following formulas \uline{in} \uline{$\unsure{\mathrm{Gl}|\mathrm{Sl}}(2,\widehat{\mathbb Z})$}

\inserted{$\det\alpha_0\in\pm1,\ \det f=1$,}
\[
\begin{gathered}
\text{(Case }h=1\text{)}\\
\left\{
\begin{aligned}
 &\alpha_0,f\in\mathrm{Gl}(2,\widehat{\mathbb Z})^{\pm1}
 \quad \deleted{\unsure{\text{with}}}\ \text{(arbitrary signs, we \uline{choose} them)}\\
 &\left.
 \begin{aligned}
 \xi&\overset{\mathrm{dfn}}=\alpha_0\rho(\alpha_0)^{-1}\\
 \varphi&\overset{\mathrm{dfn}}=f\sigma_\infty(f)^{-1}
 \end{aligned}
 \right|\ \in\mathrm{Sl}(2,\widehat{\mathbb Z})\\
 &\bigl[g\overset{\mathrm{dfn}}=\sigma_\infty(f)^{-1}\ell_1^\nu\rho(f)
 =(\ell_0^\nu f)^{-1}(\ell_0^\nu\varphi)\rho(\ell_0^\nu f)\bigr]\\
 &\text{\unsure{hence|therefore}}\quad
 \xi=\deleted{\unsure{\ell_0^{\nu}}}\cdot\varepsilon\ell_0^\nu\varphi
 \qquad\varepsilon\in\{\pm1\}\\
 &\left[\begin{aligned}
 &g\rho(g)\rho^2(g)=\varepsilon.1\\
 &g=\varepsilon.\alpha\rho(\alpha)^{-1}
 \qquad\alpha=(\ell_0^\nu f)^{-1}\alpha_0=f^{-1}\ell_0^{-\nu}\alpha_0
 \end{aligned}\right]
\end{aligned}
\right.
\end{gathered}
\tag{45}
\]

\margin{It now seems more practical to me to \uline{define} $g$ by $g=\alpha\rho(\alpha)^{-1}$, so as thus to have $g\rho(g)\rho^2(g)=1$, and to \unsure{place} the sign $\varepsilon=\varepsilon_2$ in the formula $g=\varepsilon\sigma_\infty(f^{-1})\ell_1^{\unsure{\nu}}\rho(f)$ --- \unsure{the sign} \illegible{1} \unsure{does not occur|only occurs} \illegible{3} \unsure{in} $\xi=\varepsilon\ell_0^\nu\varphi$.}
\margin{\uline{N.B.} \unsure{we write} \illegible{1}\quad $\varepsilon=\varepsilon_2$}

Symmetrically
\[
\begin{gathered}
\text{(Case }g=1\text{)}\\
\left\{
\begin{aligned}
 &\beta_0,f\in\mathrm{Sl}(2,\widehat{\mathbb Z})
 \qquad\text{(modulo signs, we \uline{choose} them)}\\
 &\eta\overset{\mathrm{dfn}}=\beta_0\sigma_\infty(\beta_0)^{-1}\\
 &\psi\overset{\mathrm{dfn}}=\unsure{\psi|f}\rho(f)^{-1}\\
 &h\overset{\mathrm{dfn}}=f^{-1}\ell_0^\nu\sigma_\infty(\rho(f))=
\end{aligned}
\right.
\end{gathered}
\tag{46}
\]

\figure{\begin{tikzpicture}[x=1mm,y=1mm,line width=.35pt,line cap=round]
\draw (0,.1) .. controls (.35,-.08) and (.9,-.06) .. (1.3,.04);
\draw (2.2,.22) .. controls (2.45,.34) and (2.65,.25) .. (2.9,.22);
\draw (4.1,.38) -- (5.4,.55);
\draw (6.7,.57) -- (7.45,.87);
\end{tikzpicture}}
\margin{147}

\scanpage{149}
\mspage{364}

I am going to look at the case $h=1$ (cf. \unsure{(45)}), setting
\[
\left\{\begin{aligned}
 \alpha_0&=\begin{pmatrix}a&b\\c&d\end{pmatrix}\\
 f&=\begin{pmatrix}r&s\\t&u\end{pmatrix}
\end{aligned}\right.
\tag{47}
\]
\margin{change of signs in $\alpha_0,\unsure{-f|f}$ is \unsure{inessential}}
\margin{N.B. In the ``arithmetic'' case (i.e. $f$ \unsure{comes|coming} from $\widehat{\Pi}_{0,3}$, $\alpha_0$ from $\widehat{\widetilde{\Pi}}_{0,3}=\widehat{\Pi}_{0,3}.\{1,\tau_\infty\}$, we will have $\alpha_0,f\equiv1\ (2)$ i.e. $a,d,r,u\equiv1\ (2)$, $b,c,s,t\equiv0\ (2)$}
\unsure{if|where} \deleted{\illegible{1}} $a,b,c,d,r,s,t,u\in\widehat{\mathbb Z}$, with
\[
 \boxed{ab-cd=\delta\in\{\pm1\}}
 \tag{48}
\]
\[
 \boxed{ru-st=1}.
 \tag{49}
\]

This \unsure{being set}
\[
 \xi=\alpha_0\rho(\alpha_0)^{-1}
 =\begin{pmatrix}A&B\\C&D\end{pmatrix}
 \tag{50}
\]
\[
 \varphi=f\underbrace{\sigma_\infty(f)^{-1}}_{{}^{t}f}
 =\begin{pmatrix}R&S\\T&\unsure{U|u}\end{pmatrix}
 =\inserted{\unsure{?}}
 \begin{pmatrix}r^2+s^2&rt+su\\\unsure{rt+su}&t^2+u^2\end{pmatrix}
 \tag{51}
\]
so that we will have (cf. \S\,42, formula (78)
\[
\left\{\begin{aligned}
 \delta A&=(a^2+b^2+ab)-(ac+bd+bc)
 &\delta B&=ac+bd+bc\\
 \delta C&=ac+bd+ad-(c^2+d^2+cd)
 &\delta D&=c^2+d^2+cd
\end{aligned}\right.,
\tag{52}
\]
\margin{N.B. we have $\delta=\unsure{\varepsilon_3}$}
and (\S\,42, (90) and 91)
\[
 R=r^2+s^2,\qquad S=T=rt+su,\qquad U=t^2+u^2.
 \tag{53}
\]
\deleted{\illegible{1}}
\deleted{\illegible{2}}\inserted{We will have} the formulas
\[
 C+D\unsure{+|-}B=1
 \tag{54}
\]
\[
 S=T.
 \tag{55}
\]
The essential relation linking $\alpha_0$ and $f$, i.e. linking the ``\unsure{variables}'' $a,b,c,d$ with $u,r,s,t$, is the relation
\[
 \boxed{\xi=\varepsilon\ell_0^\nu\varphi,\qquad\varepsilon\in\{\pm1\}}
 \tag{56}
\]
\margin{``loop relation''}
which is also written $\ell_0^{-\nu}\xi=\varepsilon\varphi$, i.e.
\margin{148}
